\exercisesection{3}

\begin{enumerate}
\def\labelenumi{\arabic{enumi}.}
\tightlist
\item
  \begin{enumerate}
  \def\labelenumii{(\alph{enumii})}
  \tightlist
  \item
    Let K be a commutative field of characteristic p \textgreater{} 0 and B the ring of formal power series with coefficients in K in two infinite systems of indeterminates \((\mathbf{X}_{n})\) , \((\mathbf{Y}_{n})\) with only a finite number of terms of given degree ( \(\S 2\) , Exercise 12); B is a Hausdorff local ring with the m-adic topology, m being its maximal ideal. Let b be the closed ideal of B generated by the monomials \(Y_{i}X_{j}\) for \(i \geqslant 0\) , \(0 \leqslant j \leqslant i\) ; let A be the local ring B/b, which is Hausdorff with the n-adic topology, where n = m/b is its maximal ideal. Let c be the element of A equal to the class mod. b of \(\sum_{n=1}^{\infty} X_{n}^{p^{n}}\) ; show that, for \(k \geqslant 3\) , \(n^{k} \cap (Ac) \not\subset n^{2}c\) (and a fortiori the relation
  \end{enumerate}
\end{enumerate}

\[
\mathfrak {n} ^ {2} (\mathfrak {n} ^ {2 k} \cap (\mathrm{A} c)) = \mathfrak {n} ^ {2 k + 2} \cap (\mathrm{A} c)
\]

cannot hold).

\begin{enumerate}
\def\labelenumi{(\alph{enumi})}
\setcounter{enumi}{1}
\tightlist
\item
  Let \(C'\) be the commutative ring with underlying set \(A \times A\) , with multiplication \((a, x)(a', x') = (aa', ax' + a'x)\) ; let C be the subring \(A \times (Ac)\) of \(C'\) ; C and \(C'\) are local rings, whose maximal ideals we denote by r and \(r'\) , and \(C'\) is a finitely generated C-module. Show that \(r = r' \cap C\) but that the topology induced on C by the \(r'\) -adic topology is not the r-adic topology.
\end{enumerate}

\begin{enumerate}
\def\labelenumi{\arabic{enumi}.}
\setcounter{enumi}{1}
\tightlist
\item
  Let A be a Noetherian integral domain which is not a field, K its field of fractions and m an ideal \(\neq\) A of A. The m-adic topology on A is not induced by the m-adic topology on K and the latter is not Hausdorff.
\end{enumerate}

* 3. If A is the ring of a non-discrete valuation of height 1 (Chapter VI) and m is its maximal ideal, A is not Hausdorff with the m-adic topology, the closure of \{0\} being m. *

¶ 4. Let A be a semi-local ring and r its Jacobson radical. For A to be Noetherian, it is necessary and sufficient that every ideal of A be closed under the r-adic topology and that every maximal ideal of A be finitely generated. (If these conditions hold, show first that the Hausdorff completion \(\hat{A}\) of A is Noetherian, using § 2, no. 13, Corollary to Proposition 19 and no. 10, Corollary 5 to Theorem 1. Then observe that the r-adic topology is Hausdorff on A and that there exists an increasing injection of the set of ideals of A into the set of ideals of \(\hat{A}\).)

¶ 5. Let A be a commutative Noetherian ring and m its Jacobson radical.

\begin{enumerate}
\def\labelenumi{(\alph{enumi})}
\item
  For A to be strictly linearly compact with a linear topology \(\mathcal{T}\) ( \(\S\) 2, Exercise 19), it is necessary and sufficient that A be semi-local and complete with the m-adic topology; \(\mathcal{T}\) is then necessarily identical with this latter topology. (Use \(\S\) 2, Exercise 21 (d) and 22 (a) and also \(\S\) 3, no. 3, Proposition 6.)
\item
  For A to be linearly compact with a linear topology T, it is necessary and sufficient that A be semi-local and complete with the m-adic topology and that T be finer than this latter topology. (To see that it is necessary, reduce it first to the case where A is a local ring, using Exercise 21 (d) of § 2. Consider first the case where T is the discrete topology and note that A is then linearly compact with the m-adic topology. In the general case, reduce it to the case where T is a minimal topology (§ 2, Exercise 18 (c)); show that A is then strictly linearly compact and use criterion (γ) of § 2, Exercise 19 (a) and Exercise 18 (b) of § 2.)
\end{enumerate}

\begin{enumerate}
\def\labelenumi{\arabic{enumi}.}
\setcounter{enumi}{5}
\item
  Let A be a Noetherian local integral domain, m its maximal ideal and K its field of fractions; suppose that A is complete with the m-adic topology. Consider on A a linear topology T which is finer than the m-adic topology. Show that, if \(T'\) is the linear topology on K for which a fundamental system of neighbourhoods of 0 consists of the neighbourhoods of 0 in A with the topology T, \(T'\) is compatible with the field structure on K and K is complete with the topology \(T'\) .
\item
  Let A be a commutative Noetherian ring and E an A-module. Let A be given the topology \(\mathcal{T}_{m}(\mathrm{A})\) (§ 2, Exercise 24) and let \(\mathcal{T}_{m}(\mathrm{E})\) denote the linear topology on E for which a fundamental system of neighbourhoods of 0 consists of the submodules a.E, where a runs through a fundamental system of neighbourhoods of 0 for \(\mathcal{T}_{m}(\mathrm{A})\) , consisting of ideals of A. Show that if E is finitely generated, E is a Hausdorff topological A-module (with \(\mathcal{T}_{m}(\mathrm{A})\) and \(\mathcal{T}_{m}(\mathrm{E})\) ) where every submodule is closed. A fundamental system of neighbourhoods of 0 for \(\mathcal{T}_{m}(\mathrm{E})\) then consists of the submodules F of E such that E/F is a module of finite length and, for every submodule M of E, the topology \(\mathcal{T}_{m}(\mathbf{M})\) is induced by \(\mathcal{T}_{m}(\mathbf{E})\) .
\end{enumerate}

¶ 8. Let A be a Noetherian semi-local ring, n its nilradical (which is the largest nilpotent ideal of A) and m its Jacobson radical. Show that, if A (with the m-adic topology) is such that A/n is complete, then A is complete. (Reduce it to the case where n is generated by a single element c such that \(c^{2} = 0\) ; using Exercise 9 of General Topology, Chapter III, §3, reduce it to showing that n = Ac is complete and use the fact that n is an (A/n)-module.)

¶ 9. (a) Let A be a Zariski ring, m a defining ideal of the topology on A and B a ring such that \(A \subset B \subset \hat{A}\) , which is a finitely generated A-module; show that, if mB is open in B under the topology induced by the topology on \(\hat{A}\) , of necessity B = A (use Nakayama's Lemma).

* (b) Let A be a commutative Noetherian ring and m an ideal of A. Show that, if the Hausdorff completion \(\hat{\mathbf{A}}\) of A with respect to the m-adic topology is a finitely generated A-module, A is complete with the m-adic topology (reduce it to the case where A is Hausdorff; use Chapter V, § 2, no. 1, Proposition 1 and Theorem 1 to show that A is a Zariski ring with the m-adic topology and conclude with the aid of (a)). *

\begin{enumerate}
\def\labelenumi{\arabic{enumi}.}
\setcounter{enumi}{9}
\item
  Let A be a Zariski ring, m a defining ideal of the topology on A and E a finitely generated A-module. Show that, if \(\hat{E}\) is an \(\hat{A}\) -module admitting a system of generators with r elements, then the A-module E admits a system of generators with r elements (note that E/mE and \(\hat{E}/m\hat{E}\) are isomorphic and use Chapter II, § 3, Corollary 2 to Proposition 4). Is the result valid if E is not assumed to be finitely generated (cf.~Exercise 9)? Is it valid if we only assume that A is Noetherian (take A = Z)?
\item
  Let A be the ring \(\mathbf{Z}_p\) of \(p\) -adic integers ( \(p\) prime) with the \(p\) -adic topology with which it is a complete Zariski ring. Let E be the A-module \(\mathrm{A}^{(N)}\) with the \(p\) -adic topology.
\end{enumerate}

\begin{enumerate}
\def\labelenumi{(\alph{enumi})}
\item
  Show that the completion \(\hat{\mathbf{E}}\) of \(\mathbf{E}\) is identified with the submodule of \(\mathbf{A}^{\mathbf{N}}\) consisting of the sequences \((a_{n})_{n\in \mathbf{N}}\) of elements of \(\mathbf{A}\) such that \(\lim a_n = 0\) .
\item
  Let \(e_n\) be the \(n\) -th vector of the canonical basis of \(E\) . Consider in \(E\) the submodule \(F\) generated by the vectors \(e_{2n-1}\) and the submodule \(G\) generated by the vectors \(p^{2n}e_{2n} - e_{2n-1}\) ( \(n \geqslant 1\) ). Show that the topologies induced on \(F\) and \(G\) by the \(p\) -adic topology on \(E\) are the \(p\) -adic topologies on \(F\) and \(G\) but that, in \(\hat{E}\) ,
\end{enumerate}

\[
\overline {{{\mathbf {F} + \mathbf {G}}}} \neq \overline {{{\mathbf {F}}}} + \overline {{{\mathbf {G}}}}.
\]

\begin{enumerate}
\def\labelenumi{(\alph{enumi})}
\setcounter{enumi}{2}
\item
  In \(\hat{\mathbf{E}}\) , let \(a_{r} = \sum_{n=0}^{\infty} p^{n} e_{r+2}\) ( \(r \geqslant 0\) ); let \(\mathbf{H}\) be the submodule of \(\hat{\mathbf{E}}\) generated by the \(a_{r}\) ( \(r \geqslant 0\) ). Show that on H the topology induced by the \(p\) -adic topology on \(\hat{\mathbf{E}}\) is the \(p\) -adic topology on H but that \(\overline{\mathbf{E} \cap \mathbf{H}} \neq \bar{\mathbf{E}} \cap \bar{\mathbf{H}}\) .
\item
  Let L be the submodule of E generated by the \(p^{n}e_{n}\) ; show that on L the topology induced by the p-adic topology on E is distinct from the p-adic topology on L.
\end{enumerate}

\begin{enumerate}
\def\labelenumi{\arabic{enumi}.}
\setcounter{enumi}{11}
\tightlist
\item
  \begin{enumerate}
  \def\labelenumii{(\alph{enumii})}
  \tightlist
  \item
    Let A be a commutative Noetherian ring and \(m_{1}\) , \(m_{2}\) two ideals of A contained in the Jacobson radical of A; let \(A_{i}\) be the completion of A with respect to the \(m_{i}\) -adic topology. If \(m_{1} \subset m_{2}\) , show that the identity mapping on A can be extended by continuity to an injective representation of \(A_{1}\) into \(A_{2}\) (cf.~General Topology, Chapter III, § 3, no. 5, Proposition 9).
  \end{enumerate}
\end{enumerate}

\begin{enumerate}
\def\labelenumi{(\alph{enumi})}
\setcounter{enumi}{1}
\tightlist
\item
  Let \(n = A_1 m_2\) . Show that, if \(A_1\) is canonically identified with a subring of \(A_2\) , \(A_2\) is identified with the completion of \(A_1\) with respect to the \(n\) -adic topology.
\end{enumerate}

\begin{enumerate}
\def\labelenumi{\arabic{enumi}.}
\setcounter{enumi}{12}
\tightlist
\item
  Let A be a Noetherian local ring with maximal ideal m and B a ring such that \(A \subset B \subset \hat{A}\) . Suppose that B is a Noetherian local ring with maximal ideal n (Exercise 14).
\end{enumerate}

\begin{enumerate}
\def\labelenumi{(\alph{enumi})}
\item
  Show that \(\mathfrak{n} = \mathbf{B} \cap \hat{\mathfrak{m}}\) and \(\mathbf{B} = \mathbf{A} + \mathfrak{n}\) . If further \(\mathfrak{n}^2 = \mathbf{B} \cap \hat{\mathfrak{m}}^2\) , then \(\mathfrak{n} = \mathbf{B}\mathfrak{m}\) (note that then \(\mathfrak{n} = \mathbf{B}\mathfrak{m} + \mathfrak{n}^2\) ).
\item
  Let K be a commutative field, C the polynomial ring K{[}X{]}, p the prime ideal CX, A the (Noetherian) local ring \(C_{p}\) and m its maximal ideal; the completion \(\hat{A}\) is identified with the ring of formal power series K{[}{[}X{]}{]} (no. 4, Proposition 8). Let \(u = u(X)\) be an element of \(\hat{A}\) which is transcendental over the field of rational fractions K(X) (cf.~Algebra, Chapter V, §5, Exercise 13 or Functions of a Real Variable, Chapter III, §1, Exercise 14 (a)) with no constant term; let B be the subring of \(\hat{A}\) consisting of the quotients
\end{enumerate}

\[
\mathrm{P} (\mathrm{X}, u (\mathrm{X})) / \mathrm{Q} (\mathrm{X}, u (\mathrm{X})),
\]

where P and Q are polynomials in K{[}X, Y{]} such that Q(0, 0) ≠ 0. Show that B is a Noetherian local ring containing A, whose maximal ideal n is generated by X and u(X); but on B the n-adic topology is strictly coarser than the m-adic topology.

\begin{enumerate}
\def\labelenumi{(\alph{enumi})}
\setcounter{enumi}{2}
\tightlist
\item
  With the same definitions as in (b), let \(\mathbf{B}'\) be the subring of \(\mathbf{B}\) generated by \(\mathbf{A}\) and \(u(\mathbf{X})\) ; show that \(\mathbf{B}'\) is not a local ring (note that \(\mathbf{B}' \cap \mathfrak{n}\) is a maximal ideal of \(\mathbf{B}'\) but that there are elements of \(\mathbf{B}'\) not invertible in \(\mathbf{B}'\) and not belonging to \(\mathbf{B}' \cap \mathfrak{n}\) ).
\end{enumerate}

¶ 14. Let K be a commutative field, C the ring K{[}X, Y{]}, p the maximal ideal CX + CY of C, A the local ring C \(_{p}\) , m its maximal ideal and \(\hat{A}\) the completion of A, which is identified with the ring of formal power series K{[}{[}X, Y{]}{]} (no. 4, Proposition 8). Let B be the subring of \(\hat{A}\) consisting of the formal power series of the type Xf(X, Y) + (P(Y)/Q(Y)), where \(f \in \hat{A}\) and P and Q are two polynomials of K{[}Y{]} such that Q(0) ≠ 0.

\begin{enumerate}
\def\labelenumi{(\alph{enumi})}
\item
  Show that \(\mathbf{A} \subset \mathbf{B}\) , that \(\mathbf{B}\) is a local ring whose maximal ideal \(\mathfrak{n}\) is equal to \(\mathbf{B}\mathfrak{m}\) and that \(\mathfrak{n}^k = \mathbf{B} \cap \hat{\mathfrak{m}}^k\) for every integer \(k \geqslant 1\) ; \(\mathbf{B}\) is therefore everywhere dense in \(\hat{\mathbf{A}}\) and the topology induced on \(\mathbf{B}\) by the \(\hat{\mathfrak{m}}\) -adic topology on \(\hat{\mathbf{A}}\) is the \(\mathfrak{n}\) -adic topology.
\item
  Show that in B the ideal b generated by all the elements of the form \(\mathbf{X}f(\mathbf{Y})\) , where \(f(\mathbf{Y}) \in \mathbf{K}[[\mathbf{Y}]]\) , is not finitely generated (cf.~Algebra, Chapter V, § 5, Exercise 13) and therefore B is not Noetherian, although B/n and \(\mathrm{gr}(\mathbf{B}) = \mathrm{gr}(\hat{\mathbf{A}})\) are Noetherian and the ideal n is finitely generated; \(b = B \cap \hat{A}X\) and b is the closure in B of the principal ideal (not closed) BX; finally B is not a flat A-module and \(\hat{B} = \hat{A}\) is not a flat B-module (use Chapter I, § 3, no. 5, Proposition 9).
\item
  Let \(f(\mathbf{Y})\) be an invertible formal power series in \(\mathbf{K}[[\mathbf{Y}]]\) , which is not an element of \(\mathbf{K}(\mathbf{Y})\) . If \(c\) is the ideal of \(\mathbf{B}\) generated by \(\mathbf{X}\) and \(\mathbf{X}f(\mathbf{Y})\) , show that on \(c\) the \(n\) -adic topology is strictly finer than the topology induced on \(c\) by the \(n\) -adic topology on \(\mathbf{B}\) . Show that the canonical mapping \(\hat{\mathbf{B}} \otimes_{\mathbf{B}} c \to \hat{c}\) (where \(\hat{c}\) is the completion of \(c\) with respect to the \(n\) -adic topology) is not injective (consider the images of \(f(\mathbf{Y}) \otimes \mathbf{X}\) and \(1 \otimes \mathbf{X}f(\mathbf{Y})\) ; to show that these two elements are distinct in \(\hat{\mathbf{B}} \otimes_{\mathbf{B}} c\) , consider this tensor product as a quotient of the tensor product \(\hat{\mathbf{B}} \otimes_{\mathbf{K}(\mathbf{Y})} c\) ).
\item
  Let \(f_{1}(\mathbf{Y}), f_{2}(\mathbf{Y})\) be two invertible formal power series in \(\mathbf{K}[[\mathbf{Y}]]\) such that \(1, f_{1}\) and \(f_{2}\) are linearly independent over \(\mathbf{K}(\mathbf{Y})\) . Let \(c_{1}, c_{2}\) be the principal ideals of \(\mathbf{B}\) generated respectively by \(\mathbf{X}f_{1}(\mathbf{Y})\) and \(\mathbf{X}f_{2}(\mathbf{Y})\) ; the n-adic topologies on \(c_{1}\) and \(c_{2}\) are respectively identical with the topologies induced by that on \(\mathbf{B}\) and the closures of these ideals in \(\hat{\mathbf{A}} = \hat{\mathbf{B}}\) are both identical with the principal ideal \(\hat{\mathbf{A}}\mathbf{X}\) of \(\hat{\mathbf{A}}\) . Deduce that the closure of \(c_{1} \cap c_{2}\) in \(\hat{\mathbf{A}}\) is not equal to \(\bar{c}_{1} \cap \bar{c}_{2}\) and that the closure of \(c_{1}: c_{2}\) is not equal to \(\bar{c}_{1}: \bar{c}_{2}\) .
\end{enumerate}

\begin{enumerate}
\def\labelenumi{\arabic{enumi}.}
\setcounter{enumi}{14}
\tightlist
\item
  \begin{enumerate}
  \def\labelenumii{(\alph{enumii})}
  \tightlist
  \item
    Let A be a Noetherian integral domain, m an ideal of A and \(\hat{A}\) the Hausdorff completion of A with respect to the m-adic topology. Show that, if M is a torsion-free A-module, then, for every element b which is not a divisor of 0 in \(\hat{A}\) , the homothety with ratio b in the \(\hat{A}\) -module \(\hat{A} \otimes_{A} M\) is injective. (Reduce it to the case where M is finitely generated; then there exists a maximal
  \end{enumerate}
\end{enumerate}

 free system \((m_j)_{1\leqslant j\leqslant n}\) in M and \(a\in \mathbf{A}\) such that for all \(m\in \mathbf{M}\) , \(am = \sum_{j}a_{j}m_{j}\) where \(a_{j}\in \mathbf{A}\) ; every \(x\in \hat{\mathbf{A}}\otimes_{\mathbf{A}}\mathbf{M}\) therefore satisfies \(ax = \sum_{j}b_{j}\otimes m_{j}\) where \(b_{j}\in \hat{\mathbf{A}}\) ; then use the flatness of the A-module \(\hat{\mathbf{A}}\) .

\begin{enumerate}
\def\labelenumi{(\alph{enumi})}
\setcounter{enumi}{1}
\tightlist
\item
  Let K be an algebraically closed field of characteristic 0, B the polynomial ring K{[}X, Y{]} and P(X, Y) = X(X² + Y²) + (X² - Y²). Show that the ideal BP is prime in B; consider the quotient ring A = B/BP which is a Noetherian integral domain. Let m be the maximal ideal of A, the canonical image of the maximal ideal n = BX + BY of B. Show that the Hausdorff completion \(\hat{A}\) of A with respect to the m-adic topology is not an integral domain (observe that in the ring of formal power series K{[}{[}X, Y{]}{]}, P decomposes into a product of two formal power series (``double point of an irreducible cubic'')). 
\end{enumerate}

\begin{enumerate}
\def\labelenumi{\arabic{enumi}.}
\setcounter{enumi}{15}
\tightlist
\item
  Let A be a commutative Noetherian ring, m an ideal of A and E a finitely generated A-module with the m-adic topology. For every infinite set I, let \(E_{I}\) denote the set of families \((x_{i})_{i\in I}\) of elements of E such that lim \(x_{i}=0\) with respect to the filter of complements of finite subsets of I; it is a submodule of the product A-module \(E^{I}\) .
\end{enumerate}

\begin{enumerate}
\def\labelenumi{(\alph{enumi})}
\item
  Show that, if \(0 \rightarrow E' \rightarrow E \rightarrow E'' \rightarrow 0\) is an exact sequence of finitely generated A-modules, the corresponding sequence \(0 \rightarrow E_{1}' \rightarrow E_{1} \rightarrow E_{1}'' \rightarrow 0\) is exact.
\item
  Define a canonical A-homomorphism \(A_{I} \otimes_{A} E \to E_{1}\) for every A-module E and show that it is an isomorphism (first verify this when E is free and finitely generated and then use (a)).
\item
  Deduce from (b) that \(A_{I}\) is a faithfully flat A-module.
\end{enumerate}

¶ 17. (a) Let K be a commutative field, A = K{[}{[}X₁, \ldots, Xₙ{]}{]} the ring of formal power series in n indeterminates with coefficients in K and V a vector space over K. With the notation of § 2, no. 6, Example 1, let V{[}{[}X₁, \ldots, Xₙ{]}{]} denote the vector space Vⁿ with the A-module structure defined by

\[
\left(\sum_ {\alpha} c _ {\alpha} \mathbf {X} ^ {\alpha}\right) (v _ {\alpha}) = (w _ {\alpha}),
\]

where \(w_{\alpha} = \sum_{\beta + \gamma = \alpha} c_{\beta} v_{\gamma}\) . Show that, if V admits a basis with I as indexing set, the A-module \(V[[X_1, \ldots, X_n]]\) is isomorphic to \(A^I\) if I is finite and to the A-module \(A_1\) defined in Exercise 16 if I is infinite. Deduce that \(V[[X_1, \ldots, X_n]]\) is a flat A-module and is faithfully flat if V is not reduced to 0.

\begin{enumerate}
\def\labelenumi{(\alph{enumi})}
\setcounter{enumi}{1}
\item
  Let L be an extension of the field K. Deduce from (a) that the ring of formal power series \(L[[X_{1},\ldots,X_{n}]]\) is a faithfully flat module over the ring \(K[[X_{1},\ldots,X_{n}]]\) . If L has finite rank over K, \(L[[X_{1},\ldots,X_{n}]]\) is isomorphic to \(L\otimes_{K}K[[X_{1},\ldots,X_{n}]]\) .
\item
  If \(\mathbf{L}\) is an algebraic extension of \(\mathbf{K}\) , show that the ring \(\mathbf{L}[[\mathbf{X}_1, \ldots, \mathbf{X}_n]]\) is a faithfully flat module over the ring
\end{enumerate}

\[
\mathbf {L} \otimes_ {\mathbf {K}} \mathbf {K} [ [ \mathbf {X} _ {1}, \dots , \mathbf {X} _ {n} ] ]
\]

(consider L as the direct limit of its sub-extensions of finite rank over K and use Chapter I, § 2, no. 7, Proposition 9). Deduce that the ring

\[
\mathbf {B} = \mathbf {L} \otimes_ {\mathbf {K}} \mathbf {K} [ [ \mathbf {X} _ {1}, \dots , \mathbf {X} _ {n} ] ]
\]

is a Noetherian local ring whose completion is identified with \(\mathbf{L}[[\mathbf{X}_1,\ldots ,\mathbf{X}_n]]\) . In order that \(\hat{\mathbf{B}} = \mathbf{B}\) , it is necessary and sufficient that \(n = 0\) or \([\mathbf{L}:\mathbf{K}] < +\infty\) .

¶ 18. Let A be a local ring with maximal ideal m; an A-module M is called quasi-finite if M/mM is a vector space of finite rank over the residue field \(k = \mathbf{A} / \mathfrak{m}\) . In particular, if \(\mathbf{A}\) is an integral domain, the field of fractions \(\mathbf{K}\) of \(\mathbf{A}\) is a quasi-finite \(\mathbf{A}\) -module.

\begin{enumerate}
\def\labelenumi{(\alph{enumi})}
\item
  Show that if A is Noetherian and M is a quasi-finite A-module, its Hausdorff completion \(\hat{\mathbf{M}}\) with respect to the m-adic topology is a finitely generated \(\hat{\mathbf{A}}\) -module (use §2, no. 11, Corollary 2 to Proposition 14). In particular, if A is complete and M Hausdorff with the m-adic topology, M is a finitely generated A-module.
\item
  Let B be another local ring, n its maximal ideal, \(\phi: A \rightarrow B\) a local homomorphism and M a finitely generated B-module. Show that, if B is Noetherian and M is a quasi-finite A-module, then the m-adic and n-adic topologies on M are identical. (Note that M/mM is a B-module of finite length and deduce that, if b is the annihilator of the B-module M/mM, then \(V(b) = \{n\}\) in Spec(B) and therefore \(V(mB + b) = \{n\}\) . Conclude using Chapter II, §4, no. 3, Corollary 2 to Proposition 11.)
\item
  Under the hypotheses of (b), show that, if \(M \neq 0\) , B/b is a quasi-finite A-module (note that \(M \neq nM\) and M/nM is a vector space of finite rank over k; deduce that B/n is of finite rank over k).
\end{enumerate}

¶ 19. Let A be a commutative Noetherian ring, m an ideal of A and S a multiplicative subset of A. Let A be a given the m-adic topology.

\begin{enumerate}
\def\labelenumi{(\alph{enumi})}
\item
  Show that the ring \(\mathbf{A}\{\mathbf{S}^{-1}\}\) (§ 2, Exercise 27) is a flat A-module.
\item
  Let \(S'\) be another multiplicative subset of \(A\) contained in \(S\) . Show that \(A\{S^{-1}\}\) is a flat \((A\{S'^{-1}\})\) -module (use Exercise 27 (d) of § 2).
\item
  Show that \(\mathbf{A}\{\mathbf{S}^{-1}\}\) is a flat \(\mathbf{A}_{\{\mathbf{S}\}}\) -module (§ 2, Exercise 27 (g)).
\item
  Suppose that \(S = A - p\) , where \(p\) is an open prime ideal of \(A\) . Show that \(A\{S^{-1}\}\) is a faithfully flat \(A_{(S)}\) -module and deduce that the ring \(A_{(S)}\) is Noetherian (cf.~§ 2, Exercise 27 (g)).
\end{enumerate}

\begin{enumerate}
\def\labelenumi{\arabic{enumi}.}
\setcounter{enumi}{19}
\item
  Let A be a commutative ring and m an ideal of A; let A be given the m-adic topology. Let B be a commutative topological A-algebra ( \(\S 2\) , Exercise 28); suppose that B is a Zariski ring; let n be a defining ideal of B. Show that, if M is a finitely generated A-module with the m-adic topology, then on the B-module \(B \otimes_{A} M\) the tensor product of the topology on B and the m-adic topology on M is the n-adic topology; then the completion tensor product \((\mathbf{B} \otimes_{\mathbf{A}} \mathbf{M})^{\wedge}\) is isomorphic to \(\hat{B} \otimes_{A} M\) .
\item
  Let A be a commutative Noetherian ring, m an ideal of A and M and N two finitely generated A-modules.
\end{enumerate}

\begin{enumerate}
\def\labelenumi{(\alph{enumi})}
\item
  Suppose that M is Hausdorff with the m-adic topology. Show that in \(\operatorname{Hom}_{\mathbf{A}}(\mathbf{M}, \mathbf{N})\) with the m-adic topology the set of injective homomorphisms is open (use no. 1, Proposition 2 and the Artin-Rees Lemma).
\item
  Suppose that A is a complete Zariski ring and m a defining ideal of
\end{enumerate}

A. For every integer \(i\) , let \(\mathbf{A}_i = \mathbf{A} / \mathfrak{m}^{i+1}\) , \(\mathbf{M}_i = \mathbf{M} / \mathfrak{m}^{i+1}\mathbf{M}\) , \(\mathbf{N}_i = \mathbf{N} / \mathfrak{m}^{i+1}\mathbf{N}\) ; show that the topological A-module \(\operatorname{Hom}_{\mathbf{A}}(\mathbf{M}, \mathbf{N})\) is isomorphic to

\[
\varprojlim \operatorname{Hom} _ {\mathbf {A}} (\mathbf {M} _ {i}, \mathbf {N} _ {i}).
\]

Deduce that in \(\operatorname{Hom}_{\mathbf{A}}(\mathbf{M},\mathbf{N})\) the set of surjective homomorphisms is open.

¶ 22. (*) Let A be a ring, commutative or not; all the A-modules to be considered are left A-modules. Let P be a property such that: (α) if f: M → N is an injective A-module homomorphism and N has property P, then M has property P; (β) the direct sum of two A-modules with property P has property P.

\begin{enumerate}
\def\labelenumi{(\alph{enumi})}
\item
  Let M be an A-module. Show that the submodules \(M'\) of M such that \(M/M'\) has property P form a fundamental system of neighbourhoods of 0 for a linear topology \(\mathcal{T}_{P}(M)\) on M. Show that \(\mathcal{T}_{P}(A_{s})\) is compatible with the ring structure on A and that M with the topology \(\mathcal{T}_{P}(M)\) is a topological module over A with \(\mathcal{T}_{P}(A_{s})\) . Every A-module homomorphism \(f: M \to N\) is continuous with the topologies \(\mathcal{T}_{P}(M)\) and \(\mathcal{T}_{P}(N)\) .
\item
  Suppose that the following condition holds: \((\gamma)\) if \(N\) is a submodule of \(M\) with property \(P\) and, for every submodule \(L \neq 0\) of \(M\) , \(N \cap L \neq 0\) , then \(M\) has property \(P\) . Show that with these conditions, if \(F\) is a submodule of an A-module \(E\) , the topology \(\mathcal{T}_P(F)\) is induced by \(\mathcal{T}_P(E)\) . (Let \(F'\) be a submodule of \(F\) such that \(F/F'\) has property \(P\) ; consider a maximal element \(G\) among the submodules of \(E\) such that \(G \cap F = F'\) and show that \(E/G\) has property \(P\) .)
\end{enumerate}

¶ 23. (a) Let A be a commutative Noetherian ring and m an ideal of A. Let P\{M\} denote the following property: M is an A-module and every finitely generated submodule of M is annihilated by a power of m.

Show that conditions \((\alpha)\) , \((\beta)\) and \((\gamma)\) of Exercise 22 are fulfilled. (To prove \((\gamma)\) , reduce it to the case where M is finitely generated; for all \(a \in \mathfrak{m}\) , there exists by hypothesis \(k > 0\) such that \(a^k\mathrm{N} = 0\) ; use the fact that there exists \(r > 0\) such that \(\operatorname{Ker}(a_{\mathbf{M}}^r) \cap \operatorname{Im}(a_{\mathbf{M}}^r) = 0\) (Algebra, Chapter VIII, § 2, no. 2, Lemma 2).)

\begin{enumerate}
\def\labelenumi{(\alph{enumi})}
\setcounter{enumi}{1}
\item
  Show that, if M is a finitely generated A-module, the topology \(\mathcal{T}_{P}(\mathbf{M})\) is identical with the m-adic topology. Give an example of an A-module M for which \(\mathcal{T}_{P}(\mathbf{M})\) is strictly finer than the m-adic topology (cf.~Exercise 11).
\item
  Show that the conclusion of (a) does not extend to the case where A is a non-commutative left Noetherian ring and m is a two-sided ideal of A (consider the ring of lower triangular matrices of order 2 over a commutative field).
\end{enumerate}

¶ 24. A ring (commutative or not) A, not reduced to 0, is called a principal ideal ring if it has no divisors of zero and every left or right ideal of A is monogenous. Such a ring is left and right Noetherian.

\begin{enumerate}
\def\labelenumi{(\alph{enumi})}
\item
  Show that, for every element \(c \neq 0\) of \(\mathbf{A}\) , \(\mathbf{A} / \mathbf{A}c\) is an \(\mathbf{A}\) -module of finite length. (Observe that, if a decreasing sequence of ideals \(\mathbf{A}a_n\) of \(\mathbf{A}\) contains \(\mathbf{A}c\) , then \(c = b_na_n\) for all \(n\) and consider the right ideals \(b_na\) .)
\item
  Show that every submodule of a (left or right) free A-module is free (same argument as in Algebra, Chapter VII, § 3, Theorem 1).
\item
  In every A-module M, the set of elements of M which are not free is a submodule T of M called the torsion submodule of M (use Chapter II, §2, Exercise 14 (a)); M is called a torsion module if T = M; M is called torsion-free if T = 0.
\item
  Show that every finitely generated torsion-free A-module is free (use Chapter II, § 2, Exercise 23 (b)).
\item
  Show that every finitely generated A-module is the direct sum of a free module and a torsion module.
\item
  Let \(\mathfrak{a}\) be a two-sided ideal \(\neq 0\) of A. Show that there exists an element \(a \in \mathfrak{a}\) and an automorphism \(\sigma\) of A such that \(\mathfrak{a} = \mathrm{A}a = a\mathrm{A}\) and \(ax = \sigma(x)a\) for all \(x \in \mathrm{A}\) . (If \(b\) is \(a\) generator of the left ideal \(\mathfrak{a}\) , there exists an endomorphism \(\tau\) of A such that \(bx = \tau(x)b\) for all \(x \in \mathrm{A}\) ; show that, if \(a = ub\) is a generator of the right ideal \(\mathfrak{a}, u\) is invertible, using Algebra, Chapter VIII, § 2, Exercise 8 (b).)
\end{enumerate}

¶ 25. Let A be a principal ideal ring (Exercise 24), a a two-sided ideal ≠0 of A and a an element of a with the properties stated in Exercise 24 (f). Let P\{M\} denote the following property: M is a left A-module and every finitely generated submodule of M is annihilated by a power of a.

\begin{enumerate}
\def\labelenumi{(\alph{enumi})}
\item
   Show that the conditions \((\alpha)\) , \((\beta)\) and \((\gamma)\) of Exercise 22 are fulfilled (consider the homothety \(a_{\mathbf{M}}\) and argue as in Exercise 23 (a), observing that \(\operatorname{Ker}(a_{\mathbf{M}}^{r})\) and \(\operatorname{Im}(a_{\mathbf{M}}^{r})\) are submodules of \(\mathbf{M}\) ).
\item
  Show that every torsion A-module M (Exercise 24 (c)) is the direct sum of a submodule \(M_{a}\) which has property P and a submodule \(M_{a}^{\prime}\) such that the restriction to \(M_{a}^{\prime}\) of the homothety \(a_{M}\) is bijective (observe that, if N is a torsion A-module and \(a_{N}\) is injective, then \(a_{N}\) is bijective; for this, reduce it to the case where N is monogenous and use Exercise 24 (a) and also Algebra, Chapter VIII, § 2, no. 2, Lemma 2).
\item
  Let S be the set of elements \(s \in A\) whose canonical image in the ring A/a is invertible. Show that S is a multiplicative subset of A and that the following conditions, for any left ideal l of A, are equivalent:
\end{enumerate}

\((\alpha) \mathfrak{l} \cap \mathbb{S} \neq \varnothing; (\beta) \mathfrak{l} + \mathfrak{a} = \mathbb{A}; (\gamma) (\mathbb{A} / \mathfrak{l})_{\mathfrak{a}} = 0\) (in the notation of (b)); (δ) for all \(x \in \mathbb{A}\) , there exists \(s \in \mathbb{S}\) such that \(sx \in \mathfrak{l}\) .

Deduce that A has a (left or right) ring of fractions with respect to S (Chapter II, § 2, Exercise 22), which is a principal ideal ring and whose only non-zero two-sided ideals are generated by the canonical images of the ideals \(a^{n}\) ; moreover the canonical mapping of A to this ring of fractions is injective.

\begin{enumerate}
\def\labelenumi{(\alph{enumi})}
\setcounter{enumi}{3}
\tightlist
\item
  Suppose now that the two-sided ideal a is maximal. Show that, for every integer n \textgreater{} 0, the ring \(A/a^{n}\) is isomorphic to a matrix ring \(\mathbf{M}_{r}(\mathbf{B}_{n})\) over a completely primary ring (Algebra, Chapter VIII, § 6, Exercise 20). If \(\mathfrak{b}_n\) is the maximal ideal of \(\mathbf{B}_n\) , show that \(\mathfrak{b}_n^n = 0\) , that every (left or right) ideal of \(\mathbf{B}_n\) is of the form \(\mathfrak{b}_n^k\) and is monogenous. (Note on the one hand that
\end{enumerate}

\[
\mathbf {a} / \mathbf {a} ^ {n} = \mathbf {M} _ {r} (\mathfrak {b} _ {n})
\]

(Algebra, Chapter VIII, § 6, Exercise 5) and on the other that, for \(k < n\) , \(\mathfrak{b}_n^k / \mathfrak{b}_n^{k+1}\) is necessarily a simple \(B_n\) -module, without which \(\mathbf{M}_r(\mathfrak{b}_n^k)\) could not be a monogenous \((A / a^n)\) -module.)

Deduce that the completion \(\hat{\mathbf{A}}\) of \(\mathbf{A}\) with respect to the topology \(\mathcal{T}_{\mathbf{P}}(\mathbf{A}_s)\) is a matrix ring \(\mathbf{M}_r(\mathbf{B})\) over a ring \(\mathbf{B}\) with no divisors of zero and each of whose ideals (left or right) is a power of the same maximal two-sided ideal.

\begin{enumerate}
\def\labelenumi{\arabic{enumi}.}
\setcounter{enumi}{25}
\tightlist
\item
  Let B be a commutative ring and A a complete Noetherian semi-local subring of B. Let n be an ideal of B containing a power of the Jacobson radical of A and such that the n-adic topology on B is Hausdorff. Then show that the topology on A is induced by the n-adic topology on B (use Proposition 8 of §2, no. 7).
\end{enumerate}

¶ 27. Let A be a ring, B a commutative A-algebra which is a Zariski ring and N a finitely generated B-module.

\begin{enumerate}
\def\labelenumi{(\alph{enumi})}
\item
  Suppose that, for an ideal \(\mathfrak{I}\) of B contained in the Jacobson radical of B, the A-modules \(\mathbf{N} / \mathfrak{I}^{n + 1}\mathbf{N}\) are flat for \(n\geqslant 0\) . Show that N is a flat A-module. (If \(v\colon \mathbf{M}\to \mathbf{M}'\) is an injective homomorphism of finitely generated A-modules, it is necessary to prove that \(u = v\otimes 1:\mathbf{M}\otimes_{\mathbf{A}}\mathbf{N}\rightarrow \mathbf{M}'\otimes_{\mathbf{A}}\mathbf{N}\) is injective. Reduce it to proving that, if we write \(\mathbf{N}_n = \mathbf{N} / \mathfrak{I}^{n + 1}\mathbf{N}\) and \(u_{n} = v\otimes 1_{\mathbf{N}_{n}}\) , the homomorphism \(\lim u_n\) is injective, using the fact that the \(\mathfrak{I}\) -adic topologies on \(\mathbf{M}\otimes_{\mathbf{A}}\mathbf{N}\) and \(\mathbf{M}'\otimes_{\mathbf{A}}\mathbf{N}\) are Hausdorff.)
\item
  Let \(b\) be an element contained in the Jacobson radical of \(\mathbf{B}\) and such that the homothety with ratio \(b\) on \(\mathbf{N}\) is injective. Show that, if \(\mathbf{N} / b\mathbf{N}\) is a flat A-module, \(\mathbf{N}\) is a flat A-module (reduce it to the case in (a)).
\item
  Suppose further that A is a local ring with maximal ideal m and residue field \(k = A/m\) and that mB is contained in the Jacobson radical of B. Let P be a B-module which is a flat A-module and \(u: N \to P\) a homomorphism such that \(u \otimes 1_k: N \otimes_A k \to P \otimes_A k\) is injective. Then show that N is a flat A-module and that u is injective. (Reduce it to showing that, for every finitely generated A-module M, the homomorphism \(u \otimes 1_M: N \otimes_A M \to P \otimes_A M\) is injective; observe that the m-adic topology on \(N \otimes_A M\) is Hausdorff; then use §2, no. 8, Corollary 1 to Theorem 1, considering the commutative diagram
\end{enumerate}

\[
\begin{array}{c} \operatorname{gr} _ {0} (\mathbf {N}) \otimes_ {k} \operatorname{gr} (\mathbf {M}) \xrightarrow {\operatorname{gr} _ {0} (u) \otimes 1} \operatorname{gr} _ {0} (\mathbf {P}) \otimes_ {k} \operatorname{gr} (\mathbf {M}) \\ \Biggl \downarrow \\ \operatorname{gr} (\mathbf {N} \otimes_ {\mathbf {A}} \mathbf {M}) \xrightarrow [ \operatorname{gr} (u \otimes 1) ]{} \operatorname{gr} (\mathbf {P} \otimes_ {\mathbf {A}} \mathbf {M}) \end{array}
\]

and using the flatness of P, the vertical arrows being the canonical homomorphisms of § 2, Exercise 6.)

